\section{LLMs for OS Security}
\label{sec:os}

\subsection{Overview and taxonomy}


\begin{figure}[t]
  \centering

  % Resizebox ensures it fits exactly into one column width
\resizebox{\columnwidth}{!}{
\begin{tikzpicture}[
    font=\sffamily\bfseries,
    % Narrower containers for column fit
    container/.style={draw, fill=#1, minimum width=2.4cm, minimum height=8.5cm, thick},
    % Smaller item boxes
    itembox/.style={draw, fill=white, minimum width=2.1cm, minimum height=0.6cm, 
                    align=center, thick, font=\sffamily\scriptsize\bfseries},
    % Arrow styles
    redarrow/.style={-Stealth, red, line width=0.8pt},
    goldarrow/.style={-Stealth, orange!90!black, line width=0.8pt},
    blackarrow/.style={-Stealth, black, line width=0.8pt}
]

    % --- Column 1: LLM cap. ---
    \node[container=blue!15] (LLM) {};
    \node[anchor=north, font=\large, yshift=-0.2cm] at (LLM.north) {\textbf{LLM cap.}};
    
    \node[itembox, yshift=1.5cm] (Reasoning) at (LLM.center) {Reasoning};
    \node[itembox, yshift=-0.5cm] (Synthesis) at (LLM.center) {Synthesis};
    \node[itembox, yshift=-2.5cm] (Automation) at (LLM.center) {Automation};

    % --- The Large Block Arrow ---
    \node[single arrow, draw, minimum height=0.5cm, minimum width=0.6cm, 
          right=0.2cm of LLM, anchor=west, inner sep=1pt, thick] (BlockArrow) {};

    % --- Column 2: OS ---
    \node[container=blue!25, right=0.9cm of LLM] (OS) {};
    \node[anchor=north, font=\large, yshift=-0.2cm] at (OS.north) {\textbf{OS}};

    % Positioned carefully to avoid "OS" title overlap
    \node[itembox, yshift=2.9cm] (Kernel) at (OS.center) {Kernel (\S\ref{sec:os:kernel})};
    \node[itembox, yshift=1.8cm] (Drivers) at (OS.center) {Device\\Drivers (\S\ref{sec:os:driver})};
    \node[itembox, yshift=0.7cm] (FileSys) at (OS.center) {File Systems \\(\S\ref{sec:os:fs})};
    \node[itembox, yshift=-0.4cm] (Networking) at (OS.center) {Networking \\(\S\ref{sec:os:network})};
  
    \node[itembox, yshift=-1.5cm] (Services) at (OS.center) {System Services \\(\S\ref{sec:os:services})};
    \node[itembox, yshift=-2.6cm] (Hypervisor) at (OS.center) {Hypervisor/\\TEE (\S\ref{sec:os:tee})};

    % --- Column 3: PL tech. ---
    \node[container=gray!15, right=1.0cm of OS] (PL) {};
    \node[anchor=north, font=\large, yshift=-0.2cm] at (PL.north) {\textbf{PL tech.}};

    \node[itembox, yshift=2.0cm, minimum width=2.3cm] (ProgAnal) at (PL.center) {Program\\analysis/\\synthesis};
    \node[itembox, yshift=-0.5cm, minimum width=2.3cm] (Fuzzing) at (PL.center) {Fuzzing};
    \node[itembox, yshift=-3.0cm, minimum width=2.3cm] (Formal) at (PL.center) {Formal\\verification};

    % --- Connections ---
    % Red arrows
    \foreach \dest in {Kernel, Drivers, Networking, FileSys, Services, Hypervisor}
        \draw[redarrow] (ProgAnal.west) -- (\dest.east);

    % Gold arrows
    \foreach \dest in {Kernel, Drivers, Networking}
        \draw[goldarrow] (Fuzzing.west) -- (\dest.east);

    % Black arrows
    \foreach \dest in {Kernel, Networking}
        \draw[blackarrow] (Formal.west) -- (\dest.east);

\end{tikzpicture}
}
  \caption{The structure of this survey, showing how fundamental LLM capabilities benefit OS security and the underpinning programming language techniques.}
  \label{fig:llm-os-mindmap}
\end{figure}

% An OS concentrates security-critical complexity: large privileged codebases, deep hardware interaction, concurrency, and decades of backwards compatibility \cite{bessey2010few}. These factors make vulnerabilities as well as logic bugs hard to prevent and expensive to triage for an OS, as signals are noisy and scattered across code, specifications, and tool outputs \cite{guo2009linux}. For instance, static analysis and address sanitizers are known hard for an OS kernel~\cite{guo2009linux,serebryany2012addresssanitizer} and to produce many false alarms~\cite{guo2009linux,wen2024automatically}.
An OS concentrates security-critical complexity: large privileged codebases, deep hardware interaction, concurrency, and decades of backwards compatibility~\cite{bessey2010few}. These factors make vulnerabilities and logic bugs hard to prevent and expensive to triage, as signals are noisy and scattered across code, specifications, and tool outputs~\cite{guo2009linux,wen2024automatically}. Static-analysis reports can require substantial triage effort, whereas AddressSanitizer detects memory errors dynamically~\cite{guo2009linux,wen2024automatically,serebryany2012addresssanitizer}.

\revise{LLMs offer a complementary way to interpret and connect these artifacts. Evaluations of code-trained models demonstrate useful code-generation capabilities, motivating their use as assistants rather than as independent correctness oracles}~\cite{chen2021evaluating}.
For example, in static-analysis-based bug finding, an LLM can summarize long call chains, explain data-structure/locking assumptions, and help validate whether an alarm is plausible by cross-checking code context and constraints~\cite{wen2024automatically}.

% while the final decision should still rely on deterministic evidence (e.g., reproducing inputs, traces, and tests).

\paragraph{Taxonomy of this survey}
\revisegreen{We organize the literature by OS component and interface, and compare studies within and across components according to the security workflows they support. The component dimension identifies the interfaces and execution environments under study, while the workflow dimension identifies the LLM's task, the artifacts it produces, and how those artifacts are evaluated.} The following sections describe their security tasks, supporting tools, and evaluation evidence.
\begin{itemize}
  \item \textbf{Kernel.} The kernel is the core of an OS. Here, we cover process and memory management, scheduling, page management, and syscall subsystems. We also consider eBPF, which permits restricted user-provided programs to execute in kernel context subject to verifier and permission checks.
  \item \textbf{Device drivers.} Drivers bridge privileged kernel code and hardware through interfaces such as \texttt{ioctl}. Their analysis must account for device-specific state and the OS and hardware environment.
  \item \textbf{Networking.} The OS networking stack processes and routes packets while maintaining stateful protocol logic, with data flowing through shared structures such as \texttt{sk\_buff}.
  \item \textbf{File systems.} File systems map files and directories onto storage devices, manage metadata such as names and permissions, and provide buffering and caching. Journaling and copy-on-write designs support crash consistency and recovery.
  \item \textbf{System services.} These user-space components often run with elevated privileges and provide logging and auditing, service management, update and configuration functions, and process orchestration. In this survey, the category covers privileged OS services.
  \item \textbf{Hypervisors and TEEs.} Hypervisors virtualize hardware for guest operating systems. TEEs provide hardware-enforced isolation for code and data. We treat both as kernel-adjacent components because their interfaces shape system trust boundaries.
\end{itemize}

\paragraph{Survey organization}
%We first provide an overview of OS security challenges and motivate why LLMs are a useful ``assistant'' for connecting heterogeneous evidence (code, documentation, static-analysis warnings, sanitizer traces, and fuzzer crashes). 
Following our component-based taxonomy (Figure~\ref{fig:os-arch}), we list the representative techniques in Table~\ref{tab:llm-kernel-related}, and describe how LLMs benefit each component: kernel (\S~\ref{sec:os:kernel}), device drivers (\S~\ref{sec:os:driver}), file systems (\S~\ref{sec:os:fs}), networking (\S~\ref{sec:os:network}), TEEs (\S~\ref{sec:os:tee}) and privileged user-space system services (\S~\ref{sec:os:services}). 
For each component, we highlight the security challenges, typical techniques/defenses, artifacts, and evaluation used in that component. 
Finally, we summarize cross-cutting themes (e.g., grounding, feedback loops, and validation) and discuss open problems and future research directions.

\revise{Table~\ref{tab:representative-evidence} compares representative studies, their evaluations, and the limits of their reported evidence.}
% Sources for numerical claims (publication snapshots, not live counts):
% KernelGPT ASPLOS 2025 abstract and Table 3: https://yangchenyuan.github.io/files/ASPLOS25-KernelGPT.pdf
% SyzAgent FASE 2025, Sec. 4 and Table 1: https://arxiv.org/pdf/2503.02301
% BugLens ASE 2025, abstract and Table I: https://www.cs.ucr.edu/~zhiyunq/pub/ase25_buglens.pdf
% KNighter SOSP 2025, Sec. 5 and Table 2: https://yangchenyuan.github.io/files/SOSP25-KNighter.pdf
% MigGPT NeurIPS 2025, Sec. 6.1 and Table 5: https://proceedings.neurips.cc/paper_files/paper/2025/file/3760dbb5835bf0b771c3f83cb27ef2c0-Paper-Conference.pdf
% PathSentinel arXiv:2407.11279, Sec. VII-D: https://arxiv.org/html/2407.11279
\begin{table*}[t]
\centering
\caption{\revise{Representative LLM-assisted OS security approaches and their reported evaluations.}}
\label{tab:representative-evidence}
\begingroup
\setlength{\tabcolsep}{2pt}
\footnotesize
\begin{tabular}{@{}>{\RaggedRight\arraybackslash}p{0.105\textwidth}>{\RaggedRight\arraybackslash}p{0.14\textwidth}>{\RaggedRight\arraybackslash}p{0.15\textwidth}>{\RaggedRight\arraybackslash}p{0.15\textwidth}>{\RaggedRight\arraybackslash}p{0.17\textwidth}>{\RaggedRight\arraybackslash}p{0.18\textwidth}@{}}
\hline
\revise{Study} & \revise{OS target and task} & \revise{LLM role and validation} & \revise{Evaluation and baseline} & \revise{Reported result} & \revise{Scope and limitation} \\
\hline
\revise{KernelGPT}~\revcite{yang2025kernelgpt} &
\revise{Linux drivers and sockets; syscall fuzzing} &
\revise{Generates and repairs Syzkaller specifications using validation feedback} &
\revise{Kernel fuzzing; Syzkaller and SyzDescribe baselines} &
\revise{24 new bugs; 11 CVEs and 12 fixes reported in the ASPLOS paper} &
\revise{Findings depend on inferred interfaces and execution coverage; accepted specifications need not be complete.} \\
\hline
\revise{SyzAgent}~\revcite{syzagent25} &
\revise{Linux kernel; directed fuzzing} &
\revise{Guides Syzkaller test generation and mutation during fuzzing} &
\revise{27 target-function cases; Syzkaller baseline; three two-hour runs per case} &
\revise{18 of 27 cases had at least 10\% relative improvement under the paper's $\omega$ metric} &
\revise{Preliminary target-reachability evaluation; the hit-rate metric is not a vulnerability count.} \\
\hline
\revise{BugLens}~\revcite{li2025towards} &
\revise{Linux kernel drivers; taint-style warning triage} &
\revise{Reasons over analyzer reports, impact and source-code constraints} &
\revise{Android Linux 4.14.150 driver reports; Suture and CodeQL-SOD} &
\revise{Suture-report precision rose from 0.10 to 0.72; four real bugs previously dismissed as false positives in Suture's manual review} &
\revise{Refines existing reports and depends on the input analyzer's dataflow; reported gains are dataset-specific.} \\
\hline
\revise{KNighter}~\revcite{yang2025knighter} &
\revise{Linux kernel; static bug-pattern detection} &
\revise{Synthesizes and refines checkers from historical bug-fix patches} &
\revise{61 bug-fix patches; generated checkers applied to the kernel} &
\revise{92 new bugs, including 77 confirmed and 30 CVEs, in the SOSP paper} &
\revise{Checker construction depends on input patches; the 32.2\% false-positive rate applies to reports from the evaluated triage pipeline.} \\
\hline
\revise{MigGPT}~\revcite{dang2025miggpt} &
\revise{Linux kernel; patch migration and security backporting} &
\revise{Retrieves relevant code and generates migrated patches} &
\revise{Out-of-tree migration and FixMorph CVE-backporting benchmark with 350 instances} &
\revise{On CVE backporting, GPT-4-turbo semantic match: 85.43\% baseline vs. 91.78\% with MigGPT} &
\revise{Semantic match is not a proof of patch correctness; outcomes depend on model and target version.} \\
\hline
\revise{PathSentinel} \revcite{lee2024static} &
\revise{Android applications and core services; filesystem-access vulnerabilities} &
\revise{Generates test exploit code using findings from static and access-control analysis} &
\revise{16 path-traversal cases; LLM-generated tests with analysis context} &
\revise{Usable tests for 7 simple and 9 complex cases with appropriate context} &
\revise{The LLM assists validation; vulnerability detection is performed by program and policy analysis.} \\
\hline
\revise{LLM-CAL}~\revcite{gadey2026llmcal} &
\revise{TEE boundary design; annotate cryptographic code for TEE placement} &
\revise{Fine-tuned LLM labels security-sensitive code to support partitioning} &
\revise{4,010 C source files; held-out line-level classification} &
\revise{98.41\% F1 and 97.50\% recall for line-level labels} &
\revise{Code-label metrics do not validate the security of a resulting TEE partition.} \\
\hline
\revise{OpenAI pipeline}~\revcite{openai2026paths} &
\revise{Linux kernel; local privilege escalation} &
\revise{Finds candidates, reproduces bugs, develops exploits and proposes patches} &
\revise{Stock Linux distribution; provider report; human review} &
\revise{1,724 reproduced vulnerabilities; 948 working LPE exploits} &
\revise{Provider-reported results depend on target configuration and resources; patches require review.} \\
\hline
\end{tabular}
\endgroup
\end{table*}

\paragraph{\revise{Cross-study synthesis.}}
\revise{The eight studies intervene at different points in a security workflow. KernelGPT improves the specifications supplied to a fuzzer, while SyzAgent steers mutation toward target functions; BugLens refines reports produced by an analyzer, KNighter generates the analyzer itself, MigGPT proposes migrated patches, PathSentinel assists exploit validation for Android applications, LLM-CAL labels cryptographic code for TEE partitioning, and the OpenAI pipeline reports exploitability assessment and candidate patches. Coverage or target hit rate, report precision, newly confirmed bugs, patch semantic match, code-label metrics and exploit counts measure different outcomes and should not be ranked as if they were a shared benchmark. Improved reachability does not by itself establish a vulnerability, agreement with a reference patch does not establish security correctness, and high code-label scores do not validate a resulting TEE boundary.}

\revise{Tool feedback supplies useful checks, but each study retains a distinct evidence boundary: fuzzing requires executable interfaces and a meaningful oracle; report triage inherits the analyzer's candidate set and dataflow; patch-guided checker synthesis depends on representative historical fixes; and patch migration must preserve the target kernel's semantics. These dependencies motivate reporting baseline configuration, dataset construction, model and kernel versions, false positives, and independent reproduction or review together with success metrics.}


% A key survey observation is that OS components differ along three axes that directly shape LLM utility: (i) \,\emph{where untrusted input enters} (system calls, ioctls, packets, and on-disk images), (ii) \,\emph{what constitutes ground truth} (crash traces, coverage, proofs, differential oracles), and (iii) \,\emph{what the security objective is} (memory safety, isolation, policy compliance, availability). Consequently, effective LLM-based methods tend to be component-aware, combining natural-language reasoning with verifiable evidence and tool feedback.

\begin{figure}[t]
  \centering
  \includegraphics[width=0.8\columnwidth]{figures/arch-crop.pdf}
  \caption{Simplified OS architecture and major kernel subsystems discussed in this survey.
  Note the highlighted components are discussed in this paper.}
  \label{fig:os-arch}
\end{figure}


\subsection{\textbf{Kernel (core services, syscalls, and eBPF)}}
\label{sec:os:kernel}
The kernel is the primary trusted computing base for commodity OS security: as mentioned earlier (\S~\ref{sec:bkgd:os}) it implements core services (e.g., scheduling, memory management, namespace management, IPC) and mediates access to hardware and global resources; 
it mainly exposes the system call layer for apps, where untrusted user-space inputs cross into privileged code, often through long-lived compatibility surfaces such as ABIs, legacy syscalls, and subsystem-specific \texttt{ioctl} and \texttt{netlink}. 
As a result, kernel security work is typically organized around (i) hardening kernel code and auditing syscall-related logic, (ii) constraining privilege and ambient authority inside the kernel, and (iii) maintaining security invariants (e.g., file permissions) under concurrency and performance constraints.

A particularly salient modern subsystem is extended BPF (eBPF), which allows \textit{restricted} user-provided programs to execute in kernel context for observability (i.e., tracing kernel activities), networking, and security enforcement. 
eBPF expands the kernel attack surface beyond traditional syscalls: it introduces a verifier with complex semantic requirements, JIT compilation, helper APIs, and shared maps that connect unprivileged or semi-privileged code to kernel state. 
This creates both opportunities (programmable security mechanisms) and risks (verifier soundness bugs, JIT issues, and subtle capability/permission interactions), making eBPF a natural focus for automated reasoning and testing.

Within the above scope, LLM assistance commonly falls into four categories:
\begin{itemize}
  %\item \textbf{Formally verified kernel}: 
  \item \textbf{Formal verification of kernels}: 
  \revise{Formal verification can establish specified kernel properties under explicit modeling assumptions.}
  Yet, writing specifications as well as theorem proving is time-consuming and tedious.   
  LLMs can assist this process by helping convert informal kernel intent (e.g., API contracts, verifier invariants, lock/refcount discipline) into candidate formal specifications, and to propose proof artifacts (lemmas, tactics, proof scripts) that are accepted or rejected by a proof assistant for validation.
  
  \item \textbf{Kernel program analysis and fuzzing}:
  Kernel program analysis and fuzzing generate high-volume, noisy signals (e.g., static-analysis warnings, sanitizer traces, and fuzzer crashes) that are difficult to triage against a large codebase.
  LLMs can help by summarizing subsystem architecture and long call chains, inferring likely data/control-flow relationships near syscall handlers, and proposing hypotheses that map low-level symptoms to candidate root causes.
  For dynamic testing, LLMs can also generate syscall- and eBPF-centric harnesses, seed programs, and stateful input sequences, and iteratively refine them using tool feedback such as coverage, verifier errors, and crash reports.
  
  \item \textbf{Kernel artifact synthesis}: 
  Kernel relies on security-sensitive auxiliary artifacts to operate, e.g., security configurations, permission lists, and eBPF program templates.
  Yet these artifacts take experienced kernel developers to compose, and can lead to kernel crashes or security breaches if not handled properly.   
  LLMs may help synthesize such artifacts at a low cost -- they translate high-level security intents into kernel configuration/policy snippets, and proposing targeted instrumentation to expose invariants or collect evidence.  
  % synthesized artifacts should be treated as untrusted until validated.
  
  % generating scaffolding code (e.g., syscall exercisers, eBPF program templates, tracing scripts), translating high-level security intent into kernel configuration/policy snippets, and proposing targeted instrumentation to expose invariants or collect evidence; synthesized artifacts should be treated as untrusted until validated.
  
  \item \textbf{Kernel patching and repair}:
  Kernel is rapidly evolving software, which undergoes constant patching and fixes, e.g., for bugs or compatibility.
  Yet, kernel patching/repairing is a high-stakes, multi-step workflow --- localize the root cause, craft a minimal fix consistent with kernel conventions, and demonstrate correctness via builds, tests, and regression reproductions.
  LLMs can help automate parts of this loop by integrating existing patches, proposing candidate patches and commit messages, drafting regression tests/reproducers from crash reports, and iteratively revising the patch in response to deterministic feedback (compiler errors, test failures, sanitizer reports, and fuzz regressions).
\end{itemize}

\paragraph{Formal verification of kernels} \revise{LLM-assisted machine-checked reasoning spans both kernel-specific and transferable work. Selene directly studies project-scale proof generation for seL4; AutoRocq reports evaluations on general verification benchmarks and Linux kernel modules; other autoformalization work proposes candidate statements or proof steps for checking by a prover}~\revcite{zhang2024selene,tu2025agentic,wu2022autoformalization}. \revise{A proof checker can reject invalid proof steps, while the correctness of the formal specification and its correspondence to the intended kernel property remain separate obligations}~\revcite{wu2022autoformalization,yang2023leandojo}.

In this workflow, the LLM proposes lemmas, tactics, proof sketches, or intermediate invariants, while a proof assistant such as Coq or Rocq checks candidate proof steps. The proof-engineering burden remains a key obstacle to wider adoption.

\revise{The following studies provide general or transferable context for kernel verification.}
\revise{VeriBench is a general benchmark for generating Lean~4 verification code from Python programs.}~\revcite{miranda2025veribench}
\revise{GPass is a Coq-based neural theorem prover rather than an LLM system.}~\revcite{chen2025gpass}
\revise{A Unit Proofing Framework proposes a research agenda for code-level verification.}~\revcite{amusuo2025unit}
\revise{ABSINT-AI studies LLM-assisted abstract interpretation for JavaScript property-access errors.}~\revcite{wang2025absint}

\revise{AutoRocq reports general verification evaluations, including Linux kernel modules, while FVEL evaluates an interactive proof environment on general code-verification tasks.}~\revcite{tu2025agentic,lin2024fvel}
\revise{WhiteFox is an LLM-assisted compiler-fuzzing method and is included only as a potentially transferable testing technique.}~\revcite{yang2024whitefox}

\begin{tcolorbox}[colback=gray!8, colframe=gray!55, boxrule=0.4pt, left=1mm, right=1mm, top=0.8mm, bottom=0.8mm, title=Key findings] From a systems-security perspective, these works highlight several recurring challenges: (i) bridging the specification gap between informal kernel intent and formal statements that are both provable and security-meaningful; (ii) scaling proof search and tactic synthesis across rapidly evolving code and configurations; and (iii) maintaining trustworthy toolchains and auditable artifacts (proof scripts, traces, and extracted invariants) in the presence of probabilistic LLM outputs. \end{tcolorbox}

\paragraph{Kernel program analysis and fuzzing} 
\revise{Program analysis and fuzzing are natural entry points for LLM assistance in kernel security because kernel code contains implicit invariants that can be difficult to reconstruct and validate.}

% Recent surveys summarize how LLMs are being integrated into vulnerability detection pipelines and how they interact with classical program-analysis techniques\cite{sheng26survey,mata25survey}.

At a high level, LLM-assisted program analysis typically uses the model to (a) summarize long call chains and cross-subsystem control flow, (b) infer candidate data-structure and locking invariants, (c) triage and post-refine static-analysis alarms, and (d) connect low-level symptoms (e.g., a sanitizer report) to plausible root causes and security impact. 
For kernels, especially the Linux kernel, these capabilities are most useful when the model’s claims are grounded in explicit evidence, e.g., identifying the syscall dispatch path
(e.g., \texttt{do\_syscall\_64} $\rightarrow \cdots \rightarrow$\texttt{\_\_x64\_sys\_*} on x86-64 Linux, with intermediate dispatch functions omitted), 
the relevant task/cred objects (\texttt{task\_struct}, \texttt{cred}), and the concrete data-structure/locking context (\texttt{struct file} / \texttt{inode} / \texttt{dentry}, \texttt{sk\_buff}, RCU, spinlocks, refcounts). 
In triage, an LLM can help connect a report (KASAN/KCSAN stack trace, syzbot repro) to likely lifetime or concurrency mistakes (e.g., missing \texttt{kref\_get}/\texttt{kfree\_rcu}, wrong lock ordering). % but should be cross-checked against compiler types, CFGs, and dynamic traces \note{cite}.

\textbf{Refinement and static-analysis noise reduction.}
A recurring theme is using LLMs to reduce the analyst burden from high-volume static-analysis alarms.
\revise{LLIFT augments static analysis for Linux-kernel use-before-initialization (UBI) detection; the paper reports four previously undiscovered UBI bugs acknowledged by the Linux community}~\revcite{li2024enhancing}.
% In addition, \textbf{KNighter} explores synthesizing executable checkers from natural-language bug-pattern descriptions to scale auditing across large kernel trees\cite{yang2025knighter}.
In addition, \textbf{KNighter} synthesizes executable static-analysis checkers from historical kernel bug-fix patches, with patch-based validation and iterative refinement~\cite{yang2025knighter}.
\revise{BugStone combines LLVM-based program analysis with LLM-guided reasoning for large-scale bug discovery}~\revcite{wu2025one}.

\revise{Together, these approaches illustrate how LLMs can help interpret analysis results or synthesize checks for kernel auditing.}

\textbf{Contextual, kernel-specific detection.}
% Beyond generic alarm refinement, several works apply LLM reasoning to kernel-specific complexity, such as precise dataflow reconstruction (e.g., taint-style detection in \textbf{\textit{Towards More Accurate Static Analysis}})\cite{li2025towards} and concurrency reasoning (e.g., \textbf{\textit{When threads meet interrupts}} for deadlock detection)\cite{ye2024threads}.
\smy{Beyond generic alarm refinement, \textbf{BugLens} uses LLM-guided post-refinement for taint-style bug detection in the Linux kernel~\cite{li2025towards}.}

\revise{Guided Vulnerability Imagination (GVI) augments vulnerability-detection training data with synthesized bug patterns for its evaluated software targets.}~\revcite{yong2025gvi}

Broader dataflow-centric reasoning capabilities have also been explored in general settings (e.g., \textbf{LLMDFA}) and can be adapted to kernel-style code where pointer-rich dataflow, aliasing, and lifetimes across subsystems (VFS, networking, drivers) drive many security-relevant behaviors\cite{wang2024llmdfa}. 
In practice, adapting these methods to the kernel often requires making the analysis ``kernel-aware'' (e.g., modeling common primitives such as refcounts, RCU, and per-CPU data) and grounding outputs in concrete code references.

\revise{More general frameworks such as Fuzz4All study LLM-assisted fuzzing across programming languages and parsers.}~\revcite{xia2024fuzz4all}
% Furthermore, the reasoning capabilities of LLMs have been specialized for directed fuzzing; for instance, \textbf{RandLuzz} ~\cite{feng2025fuzzing} employs LLMs to generate high-quality seeds that satisfy complex path constraints, significantly improving the efficiency of reaching deep target code locations needed by kernel code. 
Furthermore, the reasoning capabilities of LLMs have been specialized for directed fuzzing; for instance, \textbf{RandLuzz} ~\cite{feng2025fuzzing} employs LLMs to generate high-quality seeds that satisfy complex path constraints, \smy{improving targeted exploration in its evaluated programs.} 
\revise{HLPFuzz applies LLM-based constraint solving to language-processor fuzzing}~\revcite{yang2025hybrid}. \revise{These general techniques may inform kernel fuzzing, subject to evaluation on kernel targets.}
% \textbf{Snowplow} further strengthens kernel fuzzing by learning a white-box test mutator that predicts promising argument mutations from the current/desired coverage context\cite{gong2025snowplow}.

Specification- or oracle-driven fuzzing is particularly relevant to kernels, where the key bottleneck is not input generation but deciding whether behavior is incorrect.
eBPF is emblematic: its semantics are constrained by a verifier and by non-trivial safety rules, motivating specification-based oracles for detecting misbehavior.
% \textbf{SpecCheck} is a recent and representative example of this line of work, using an explicit specification-derived oracle to flag semantic misbehavior beyond crash-based signals~\cite{lyu2025ebpf}.
% While it is currently still based on manually devised rules, such approaches align well with LLMs: a model can help draft candidate specifications, generate diverse eBPF programs that explore corner cases, and explain misbehavior reports, while the oracle enforces correctness~\cite{lyu2025ebpf}.
\smy{As non-LLM background on specification-based testing, \textbf{SpecCheck}, part of \textbf{Veritas}, checks eBPF-verifier behavior against a specification-derived oracle~\cite{lyu2025ebpf}.}

\begin{tcolorbox}[colback=gray!8, colframe=gray!55, boxrule=0.4pt, left=1mm, right=1mm, top=0.8mm, bottom=0.8mm, title=Key findings]
\revise{Across the studies reviewed here, LLMs assist with navigating subsystems, proposing hypotheses, and generating structured tests. Correctness and security conclusions still depend on task-specific evidence such as checked proofs, test oracles, traces, and reproduced crashes.}
\end{tcolorbox}

\paragraph{\revisenew{Patch-based vulnerability reproduction}}
\revisenew{K-Repro takes a kernel fix commit as input and constructs a PoC through source inspection, VM execution and debugging. On 100 KernelCTF cases, its primary configuration reproduces 56 cases in the initial run and 64 cumulatively after repeated runs. Its analysis identifies concurrency and object-lifetime constraints as recurring obstacles to reproduction}~\newrevcite{pu2026krepro}.

\revisenew{SyzHarness generates a parameterized harness from the patch and source context, preserves the prerequisite setup logic, and exposes uncertain inputs to Syzkaller mutation. Coverage feedback guides harness refinement. Its evaluation confirms 78 of 100 KernelCTF reproductions by checking failure behavior before and after the corresponding patch and matching the execution evidence to the affected code}~\newrevcite{li2026syzharness}.

\revisenew{These methods allocate the reproduction effort differently: K-Repro iterates over executable PoCs using debugging feedback, whereas SyzHarness separates setup construction from concrete input search. The latter gives the fuzzer a constrained search space while retaining LLM assistance for repairing missing setup steps. Their reported outcomes should be read with the model, search budget and reproduction criterion in Table}~\ref{tab:representative-evidence}.

\paragraph{\revise{Exploitability assessment and exploit development}}
\revise{Recent provider reports extend LLM-assisted kernel analysis to exploitability assessment.} \revisenew{OpenAI reports that its pipeline analyzed 2,308 candidate findings and dynamically reproduced 1,724 in a Linux distribution. It generated 1,169 working local privilege-escalation exploits, corresponding to 948 unique LPE vulnerabilities counted using distinct patches after variant analysis.} \revise{The pipeline generated candidate patches for human review; the reported run cost approximately 1.21 million US dollars in API credits}~\revcite{openai2026paths}.
\revise{Anthropic reports that Claude Mythos Preview constructed Linux local privilege-escalation chains involving multiple vulnerabilities, while other remotely triggerable memory vulnerabilities resisted its exploitation attempts}~\revcite{anthropic2026mythos}.
\revise{These results illustrate both offensive capability and defensive validation for prioritizing fixes. They are provider-reported observations under specific targets and resource budgets; a working exploit demonstrates impact in its tested environment, and a generated patch still requires maintainer review and validation.}

\paragraph{Kernel auxiliary artifact synthesis}
\revise{The kernel-specific synthesis work reviewed here generates auxiliary artifacts, such as static-analysis checkers and eBPF extensions}~\revcite{yang2025knighter,zheng2024kgent}.

% \textbf{Checker synthesis for kernel auditing.} KNighter exemplifies ``analysis checker synthesis'' in which an LLM derives bug patterns from historical kernel bug-fix patches and synthesizes executable static-analysis checkers, enabling scalable auditing of large kernel codebases~\cite{yang2025knighter}. Conceptually, this treats the LLM as a compiler from informal security intent to formalized detection logic; the key technical challenge is constraining the synthesized checker so that it is both sound enough to be useful and precise enough to avoid overwhelming analysts with noise ~\cite{yang2025knighter}.
\textbf{Checker synthesis for kernel auditing.} KNighter exemplifies ``analysis checker synthesis'' in which an LLM derives bug patterns from historical kernel bug-fix patches and synthesizes executable static-analysis checkers, enabling scalable auditing of large kernel codebases~\cite{yang2025knighter}. \smy{Conceptually, this turns patch-derived bug patterns into executable static-analysis checkers;} the key technical challenge is constraining the synthesized checker so that it is both sound enough to be useful and precise enough to avoid overwhelming analysts with noise ~\cite{yang2025knighter}.

% \textbf{Exploit and pathology synthesis for security assessment.}
% Beyond defensive tooling, synthesis techniques are also applied to reconstruct and generate exploits or exploit-enabling program behaviors, including data-oriented exploits that target non-control-flow data structures to bypass modern mitigations. 
% For example, \textbf{Lancet} formalizes the reconstruction of logical paths from crashes to exploit primitives \cite{dai2025lancet}. 
% More broadly, recent work demonstrates sound and efficient exploit synthesis by framing the task as a programming language synthesis problem, where the model must satisfy complex data-flow constraints to produce viable attack payloads ~\cite{bailluet2025nothing}. 
% This application underscores how synthesis can automate the evaluation of kernel resilience against sophisticated data-corruption threats.

\textbf{Kernel extension and monitoring synthesis via eBPF agents.}
The closest setting to ``kernel code synthesis'' arises with safe extension mechanisms such as eBPF ~\cite{zheng2024kgent}. Since eBPF programs can be loaded at runtime (subject to verifier constraints), they provide a constrained substrate where LLMs can synthesize kernel-resident monitoring or enforcement logic without rebuilding the kernel ~\cite{gao2025offloading}.
\revise{Kgent, initially circulated as KEN, illustrates an agentic approach in which an LLM produces and refines eBPF programs for user-specified kernel extension tasks. Its evaluation focuses on program correctness; applying generated extensions to security enforcement requires separate validation}~\revcite{zheng2024kgent}.
% Related systems (e.g., \textbf{AgentSight}) similarly use eBPF as a safe substrate to synthesize kernel-resident monitoring/observability logic \cite{zheng2025agentsight}. 
This line of work highlights the synthesis--verification interplay: eBPF's verifier and helper APIs provide guardrails, but subtle semantic gaps (kernel version differences, helper availability, and non-functional constraints such as overhead) still demand testing and careful operational controls.



\begin{tcolorbox}[colback=gray!8, colframe=gray!55, boxrule=0.4pt, left=1mm, right=1mm, top=0.8mm, bottom=0.8mm, title=Key findings] 

\revise{The evaluated synthesis systems turn bug-fix patterns or user requests into executable checkers or constrained extensions.}
\end{tcolorbox}

\paragraph{Kernel patching and repair}



\revise{Kernel live-update work studies safe application of changes to running systems, while migration and backporting studies evaluate patch adaptation across kernel versions. These tasks make the running system and target version part of the validation context}~\revcite{arnold2009ksplice,dang2025miggpt,li2025portgpt}.
% As a result, LLM-based patching for kernels is increasingly framed as an \emph{agentic} workflow in which the model iteratively proposes edits, consults tool feedback (builds, tests, sanitizers, fuzzers), and revises until a candidate patch is both correct and reviewable \cite{mathai2025crashfixer,nong2025appatch}.
\revise{A representative kernel-repair system, kAgent (formerly CrashFixer), proposes candidate patches for Linux kernel crashes and refines them using build and reproducer feedback}~\revcite{mathai2025crashfixer}.

\revisenew{Live-kBench evaluates kernel crash repair through a continuously updated bug collection and kEnv, a shared environment for building kernels and executing reproducers. Its initial Live-kBench-2512 release contains 534 bugs. It measures crash resolution through repeated execution, root-cause localization through patch overlap, and equivalence to developer fixes through an LLM judge. Separate reporting of these outcomes helps identify whether an agent removes the observed failure and how its changes relate to the developer's repair}~\newrevcite{huang2026livekbench}.

Within the emerging domain of memory-safe kernel development, \textbf{Rust-doctor} \cite{yang2025rust} addresses the unique complexities of the Rust compiler's borrow checker by automating the repair of ownership and lifetime errors through balanced training data and targeted LLM reasoning. Related work explores structured reasoning and iterative refinement for vulnerability repair (e.g., \textbf{LoopRepair}~\cite{looprepair25}), as well as approaches that explicitly decompose repair reasoning into intermediate steps (e.g., tree-of-thought style analysis in \textbf{SAN2PATCH})\cite{kim2025logs}. 
% Prompt- and interaction-centric methods such as \textbf{APPATCH} and \textbf{PatchAgent} emphasize that patch quality often depends as much on the prompting strategy and feedback integration as on raw model capacity, particularly when mimicking developer workflows (hypothesis, edit, test, and explain)\cite{nong2025appatch,xue2025pagent}.
\smy{\textbf{APPATCH} studies LLM-based vulnerability repair~\cite{nong2025appatch}, while \textbf{PAGENT} combines repository-level static analysis with LLM-based type inference to repair type-related errors in patches generated by software-engineering agents~\cite{xue2025pagent}.}

Beyond generating a single fix, several kernel-specific tasks involve reasoning over patch \emph{ecosystems} \cite{dang2025miggpt}. 
Migration and backporting are critical in practice because security fixes must be propagated across versions and vendor trees \cite{dang2025miggpt,li2025portgpt}. 
\textbf{MigGPT} targets migration of out-of-tree Linux kernel patches across versions\cite{dang2025miggpt}, while \textbf{PortGPT} focuses on automated backporting of kernel patches to older stable branches\cite{li2025portgpt}. 
\revise{MigGPT also reports a CVE patch-backporting evaluation, distinct from its broader patch-migration task}~\revcite{dang2025miggpt}.
Finally, ``patch intelligence'' tasks such as stable-patch classification can be used to triage which commits should be backported or prioritized; 
\textbf{PatchScope} exemplifies this direction by providing fine-grained classification of Linux stable patches using LLM enhancement\cite{liu2025patchscope}.


% Finally, LLMs facilitate the assessment of exploitability across complex environments. To manage long-term maintenance, \textbf{AEM} leverages LLMs for cross-version exploitability assessment, determining if a root cause persists across kernel releases despite significant code drift\cite{jiang2023aem}. 
% By anchoring probabilistic model outputs in deterministic evidence, these approaches transform raw crash data into structured security intelligence.


\begin{tcolorbox}[colback=gray!8, colframe=gray!55, boxrule=0.4pt, left=1mm, right=1mm, top=0.8mm, bottom=0.8mm, title=Key findings] 
\revisenew{Kernel repair covers both generating fixes from crash evidence and adapting existing patches across versions. Reproducer execution and comparison with developer fixes assess different aspects of a repair; Live-kBench makes these outcomes and their variation over time available in a shared evaluation environment}~\newrevcite{huang2026livekbench}.
% because they act as semantic integrators over scattered evidence—turning crash reports and sanitizer traces into coherent hypotheses, proposing localized edits, and iteratively improving patch candidates using tool feedback (builds, tests, sanitizers, fuzzers) until the result is reviewable. 
% This “agentic” generate–validate loop aligns well with kernel workflows and helps navigate patch ecosystems (e.g., migration/backporting and stable-patch triage) . 
\end{tcolorbox}


\subsection{\textbf{Device Drivers}}
\label{sec:os:driver}

Device drivers sit at the boundary between kernel software and physical devices.
They mediate stateful interactions under elevated privileges~\cite{swift2003improving}.
\revise{Driver analysis must account for device-specific state, kernel interfaces, and models of the OS and hardware environment; prior work studies these modeling requirements and ioctl-mediated driver testing}~\revcite{ball2006thorough,renzelmann2012symdrive,corina2017difuze}.
% A fundamental challenge underlying these difficulties is the specification--implementation gap.
\revise{Static analysis of drivers requires models of the operating system and device environment; building and maintaining these models is challenging across diverse hardware}~\revcite{ball2006thorough}.
\revise{SymDrive documents path-exploration and environment-modeling challenges for driver analysis, including complex ioctl inputs and interrupt behavior}~\revcite{renzelmann2012symdrive}. \revise{DiFuze takes a complementary approach by using driver-interface information to guide fuzzing}~\revcite{corina2017difuze}.

Recent research explores the use of large language models (LLMs) to augment traditional driver analysis techniques, positioning LLMs as a heuristic bridge when specifications are missing or incomplete.
\revise{These approaches use semantic regularities in source code, documentation, and development artifacts to reduce modeling gaps.}
% This makes careful example--claim alignment especially important in a survey and suggests the area is still emerging.
\revise{The following discussion focuses on driver-facing surfaces, such as \texttt{ioctl}-mediated entry points, and auditing workflows evaluated on kernel device drivers.}

\textbf{\revise{Fuzzing with incomplete specifications.}}
Coverage-guided fuzzers struggle on driver attack surfaces when valid inputs require non-trivial structure and cross-field dependencies.
\revise{KernelGPT addresses part of this challenge.}
It improves Linux kernel fuzzing by using LLMs to automatically synthesize and refine Syzkaller syscall specifications (i.e., syzlang) from implementation artifacts, and to repair these specifications based on validation feedback~\revcite{yang2025kernelgpt}.
This matters for driver security because many driver code paths are exercised through syscall-mediated entry points, especially \texttt{ioctl}, where better specifications can help fuzzers generate well-formed call sequences that reach deeper driver handlers.
% We view this line of work as strengthening the specification layer that fuzzers rely on, not as complete recovery of all driver-specific \texttt{ioctl} semantics.

\paragraph{Literature-based workflow example: KernelGPT}
\revise{KernelGPT's published workflow illustrates how generated specifications are validated and used for fuzzing; this literature-based example is not an independent empirical validation of the survey's taxonomy}~\cite{yang2025kernelgpt}.
\begin{myenumerate}
  \item \textbf{Inputs.} \revise{Kernel implementation artifacts provide evidence about syscall interfaces and the structures needed to exercise kernel and driver code.}
  \item \textbf{LLM-generated artifact.} \revise{The model proposes Syzkaller syscall specifications in syzlang, translating implementation information into inputs usable by the fuzzer.}
  \item \textbf{Validation and refinement.} \revise{Specification-validation feedback guides subsequent corrections, keeping the model's proposals within a tool-checked refinement loop.}
  \item \textbf{Downstream use.} \revise{The resulting specifications support syscall-based fuzzing, including driver-facing entry points. The workflow connects source-level interface information to executable testing artifacts.}
\end{myenumerate}
\revise{The generated specifications and validation feedback should remain inspectable: accepted specifications may omit interface semantics, and failures require reproduction and root-cause analysis before a vulnerability is confirmed.}

\textbf{Static analysis on kernel device drivers.}
LLMs have also been used as post-refinement components that reduce analyst burden without replacing conventional analyzers.
\textbf{BugLens} evaluates on Linux kernel device drivers by taking warnings produced by prior static analyzers and guiding an LLM through structured reasoning to assess security impact and validate constraints against the surrounding code context.
% In this setting, LLMs act as a semantic triage layer that filters spurious reports and highlights actionable cases~\cite{buglens25}.
In this setting, LLMs act as a semantic triage layer that filters spurious reports and highlights actionable cases~\revcite{li2025towards}.
This is useful for driver auditing, where benign low-level idioms and performance-driven patterns can confuse scalable analyses.

% Despite promising results, open challenges remain regarding robustness, scalability, and the trustworthiness of LLM-assisted reasoning in high-privilege kernel code.



% \zmy{After reviewing the relevant literature, only KernelGPT and BugLens address device drivers. Thus, I retain only these two works and highlight that LLM-based research for device driver security remains largely unexplored.}

\begin{tcolorbox}[colback=gray!8, colframe=gray!55, boxrule=0.4pt, left=1mm, right=1mm, top=0.8mm, bottom=0.8mm, title=Key findings]
\revise{The driver-facing studies discussed here use LLMs to synthesize fuzzing specifications for driver-reaching entry points or to refine static-analysis results.}
% Across these categories, LLMs primarily serve as augmentation mechanisms rather than standalone solutions, complementing traditional analyses by compensating for incomplete specifications and high-noise auditing pipelines.
\revise{Additional evaluation on diverse drivers and hardware remains an open research direction.}
% Research on LLM for Device Drivers Security is scarce, making it a field that urgently requires further investigation.
\end{tcolorbox}





\subsection{\textbf{File systems}}
\label{sec:os:fs}

File systems manage on-disk data formats for both efficiency and security. 
\revise{File-system implementations have reported security vulnerabilities across multiple components. Cai et al. analyze 157 CVE cases and characterize their attack procedures, root causes, consequences, and mitigations}~\revcite{cai2019understanding}.
% improving the data locality for various workloads and scenarios, e.g., F2FS specializes hot/cold regions for flash longevity\cite{f2fs}. 
\revise{File-system reasoning also involves invariants across on-disk formats, journaling and recovery logic, and interactions among VFS, block-layer, and page-cache code. These relationships motivate possible LLM assistance.}

\revise{Some recent work studies LLMs for file-system efficiency. STELLAR tunes parameters of parallel file systems for HPC workloads, a task its authors describe as complex and manpower-intensive}~\revcite{egersdoerfer2025stellar}.
% STELLAR achieves human-like efficiency by functioning as an agentic workflow that can read software manuals, analyze I/O trace logs, and execute trial runs on real hardware to refine its strategies.
There are only limited work which directly applies LLMs to file systems for better security, which we next describe in detail.


\revise{The most closely related work is PathSentinel, which analyzes filesystem-access vulnerabilities in Android applications and core services, including path traversal, hijacking, and luring attacks. Static program analysis and access-control analysis identify candidate issues; the LLM is used to generate exploit tests for validation, rather than to detect vulnerabilities or to analyze the filesystem implementation itself}~\revcite{lee2024static}.
% Rigorous testing on Samsung and OnePlus devices running Android 12 and 14 uncovered 51 previously unknown vulnerabilities with an exceptionally high accuracy rate. 
\revise{This work illustrates how program analysis and LLM-assisted exploit generation can support validation of the targeted Android filesystem vulnerabilities.}

\revise{Egersdoerfer et al. investigate ChatGPT-based anomaly detection on parallel-file-system logs in a preliminary study}~\cite{egersdoerfer2023early}.
% Early exploration of using chatgpt for log-based anomaly detection on parallel file systems logs, HPDC'23
\revise{The approach uses summary-based memory to retain context across sequential log windows.}
\revise{It prompts the model for both a status prediction and an explanation, giving administrators additional information to inspect alongside the predicted anomaly.}
\revise{The preliminary evaluation reports promising results against the tested baselines.}

\revise{SysGPT is a fine-tuned model for serial-performance optimization; the paper includes a file and storage systems case study}~\revcite{park2025principles}.
% By transforming performance optimization from an intuition-driven task into a structured process, this work helps developers identify and implement targeted improvements for the sequential bottlenecks that limit overall system speed as well as giving security advice of file systems.
By transforming performance optimization from an intuition-driven task into a structured process, this work helps developers identify and implement targeted improvements for the sequential bottlenecks that limit overall system speed~\cite{park2025principles}.

\revise{SysSpec uses LLMs to turn high-level designs into file-system code and adapts specification principles from formal methods. Its evaluation concerns functional correctness and system evolution, not security properties; applying its approach to security invariants remains an open direction}~\revcite{liu2025sharpen}.

\begin{tcolorbox}[colback=gray!8, colframe=gray!55, boxrule=0.4pt, left=1mm, right=1mm, top=0.8mm, bottom=0.8mm, title=Key findings]
\revise{The file-system work discussed here focuses more on efficiency and operations than on direct security tasks. PathSentinel uses an LLM to assist validation of Android application filesystem-access findings.}
\end{tcolorbox}



\subsection{\textbf{Networking}}
\label{sec:os:network}

\revisenew{An OS network stack processes packets through kernel networking and transport protocols and their packet-buffer abstractions. This section covers testing of Linux socket interfaces and analysis of protocol implementations. KnitFuzz directly exercises kernel networking, while studies such as ChatAFL evaluate application-layer protocols}~\newrevcite{zhang2026knitfuzz,meng24}.

\revise{For stateful protocols, valid message sequences can reach implementation states that malformed inputs do not; application-layer studies such as ChatAFL use state-aware generation to explore these behaviors}~\revcite{meng24}.
\revisenew{LLMs contribute executable socket tests as well as models and grammars derived from protocol specifications.}

\paragraph{RFC-Based Specification Formalization}
Similarly, network stack can also use formal verification for better security;
differently, the formal verification of network stack  must bridge the semantic gap between implementations and external natural-language standards (e.g., RFCs). 
Recent work explores using LLMs to parse RFC prose, extract requirements, and translate them into verification-ready representations.

For instance, \textbf{Mao et al.} targets RFC-derived requirements by converting natural-language specifications into models that can be consumed by protocol verification workflows to check whether intended security properties hold under the modeled assumptions~\cite{mao2025llm}.
By formalizing natural-language requirements into logic-oriented models, such approaches can support more rigorous assessment of whether a protocol design or implementation aligns with its intended security guarantees.

\revise{PARVAL uses a multi-agent framework to extract format specifications from C source and RFC documents, then checks parser inconsistencies in an FRRouting BFD-parser case study}~\revcite{zheng25}.
Similarly, \textbf{CryptoFormalEval} benchmarks the capability of LLMs to translate cryptographic protocol descriptions into formal syntax, such as Tamarin Prover, allowing for automated discovery of logical flaws through interaction with the formal toolchain~\cite{curaba24}.

\paragraph{Protocol State Inference and Fuzzing.}
\revisenew{KnitFuzz targets state dependencies in Linux socket system calls by generating C seed programs with an LLM. It combines selective memory instrumentation for mutation with a state-aware remote endpoint for network interactions. The published abstract reports two kernel bugs, including one previously unknown vulnerability, across its evaluated Linux kernels. KernelGPT supplies a fuzzer with interface descriptions; KnitFuzz supplies executable programs that establish socket-call context}~\newrevcite{zhang2026knitfuzz,yang2025kernelgpt}.

In the domain of fuzzing, LLMs also enhance coverage by helping infer protocol structure (e.g., grammars or state machines) and generating structured message sequences. 
This capability is important because traditional mutation-based fuzzers often fail to synthesize protocol-valid packets required to traverse deep states without prior knowledge of the protocol's grammar and transition constraints. 

\revisenew{Application-layer protocol testing addresses message structure and session state.}
\revise{As a non-LLM comparison, HDHunter uses gray-box generation of structured request/response sequences to expose HTTP desynchronization flaws in parsers and intermediaries}~\revcite{mu2025silent}.
\textbf{ChatAFL} demonstrates this by extracting machine-readable grammars from RFCs to guide state exploration in protocols like RTSP and FTP, and reporting improvements over baseline fuzzing strategies~\cite{meng24}. 
To address undocumented or private protocols, \textbf{LLM-Boofuzz} parses real traffic captures to synthesize executable fuzzing scripts and employs an agentic workflow to iteratively optimize test cases based on runtime feedback~\cite{wang25boofuzz}. 
Furthermore, \textbf{mGPTFuzz} demonstrates scalability by processing over a thousand pages of the Matter IoT specification to construct finite state machines, enabling stateful testing that can reveal non-crash logical vulnerabilities~\cite{ma24}.


\begin{comment}
\note{no notable works. Thus directly omit.}
\paragraph{Intent-Driven Configuration and Synthesis}
Synthesis in networking primarily manifests as Intent-Based Networking (IBN), where systems translate high-level natural-language operational goals into low-level, vendor-specific configurations and policies. 
By allowing operators to specify goals without mastering heterogeneous command-line syntaxes, this abstraction can reduce human-induced misconfigurations which is a common contributor to network-level security incidents.
In this setting, LLMs can assist by producing candidate configurations, policy templates, and human-readable rationales, but correctness should be ensured by downstream validation (e.g., configuration analyzers, simulators, or change-control pipelines) rather than by the LLM output alone.
\end{comment}


% Additionally, system-level primitives such as \textbf{SparSamp} for provably secure steganography demonstrate that LLM-adjacent generative mechanisms can be paired with formally grounded security goals, illustrating a path toward defensively motivated designs beyond bug finding~\cite{zheng25}.




\begin{tcolorbox}[colback=gray!8, colframe=gray!55, boxrule=0.4pt, left=1mm, right=1mm, top=0.8mm, bottom=0.8mm, title=Key findings]
\revisenew{The generated artifact follows the tested interface: KnitFuzz constructs C programs for kernel socket state, while application-layer fuzzers generate protocol grammars and message sequences. Both need to establish valid interaction states before exploring deeper behavior.}
\end{tcolorbox}



%\subsection{\textbf{Virtualization and trusted execution (hypervisors and TEEs): exploits and system support for confidential LLM inference)}}
\subsection{\textbf{Hypervisors and TEEs}}
% ): exploits and system support for confidential LLM inference)}}
\label{sec:os:tee}

% While Large Language Models (LLMs) have demonstrated significant proficiency in automating Kernel security tasks—such as synthesizing system call harnesses and patching memory safety vulnerabilities—their application to Hypervisors and TEEs represents a shift toward hardware-software co-design security. 
\revise{The security of hypervisors and TEE software, such as Xen}~\cite{barham2003xen} \revise{and OP-TEE}~\cite{optee}, \revise{is important because these components enforce privileged isolation boundaries. Compromises can expose guest or protected computations, depending on the affected component and threat model.}

\revise{For hypervisors, Nyx combines Intel Processor Trace coverage feedback with fast snapshot restoration, while HYPERPILL explores hardware-virtualization interfaces to generate and inject test inputs}~\revcite{schumilo2021nyx,bulekov2024hyperpill}.
\revise{For TEE interfaces, TeeDFuzzer tests interactions between the Rich Execution Environment and the TEE}~\revcite{wen2025teedfuzzer}.
\revise{TEE-Rex analyzes SGX enclave code, whereas BootStomp and Firmalice focus on secure-boot or bootloader firmware}~\revcite{Cloosters2020TeeRexDA,redini2017bootstomp,shoshitaishvili2015firmalice}.

\revise{LLM-CAL uses fine-tuned LLMs to identify cryptographic code regions for placement inside TEEs, supporting code partitioning and TCB reduction. Its evaluation measures code-label accuracy, not the security of a resulting TEE partition.}~\revcite{gadey2026llmcal}
\revise{Under the search approach described in Section~\ref{sec:methodology}, we identified LLM-CAL as work on TEE code partitioning, but did not identify evaluations in which LLMs detect vulnerabilities in hypervisor or TEE implementations or validate their security properties. The recorded queries included TEE and hardware-isolation terminology alongside LLM terms, but did not document comprehensive searches for neighboring terms such as VMMs, hypercalls, enclave interfaces, and secure monitors; recall may therefore be limited.}
\paragraph{Adjacent work: protecting LLMs with TEEs}
\revise{The following systems use isolation mechanisms to protect LLM execution. They represent a separate direction from using LLMs to help design TEE boundaries or secure OS components.}

% Unlike the Kernel, where the challenge lies in managing vast codebases and complex business logic, the security of TEEs and Hypervisors hinges on strict adherence to low-level architectural specifications and cross-privilege boundary isolation. 
% LLM-assisted frameworks in this domain act as 'semantic bridges,' translating unstructured hardware manuals and RFCs into formal models or state-aware fuzzing seeds. 
% This enables the discovery of deep-seated logical flaws in hypercalls and enclave entries that traditional coverage-guided tools, which lack understanding of hardware-dependent execution constraints, often fail to reach.

%\paragraph{Exploit and assessment context}
%For LLM deployments, hypervisors and TEEs defend \emph{confidentiality and integrity} of model weights, prompts, and intermediate state against an untrusted host OS or co-tenant software: hypervisors provide VM-level isolation and controlled virtual devices, while TEEs aim to isolate code and data even from a privileged host. 
%However, they do not eliminate OS-style risks: the VM/host and host/TEE boundaries become syscall-like attack surfaces, and microarchitectural side channels can still leak input-dependent access patterns. 
%Beyond TEE-specific leakage, LLM-specific failure modes also matter when assessing ``confidential inference'': positional/format biases can systematically skew outputs (e.g., multiple-choice positional bias as studied in \textbf{Anchored answers}~\note{cite}), and privacy properties are increasingly treated as first-class security objectives (e.g., data removal guarantees in \textbf{LLMs are Privacy Erasable}~\cite{ye2025llms}). 
%Together, these lines motivate evaluating TEE/virtualization deployments not only for enclave escape, VM escape, or key theft, but also for what residual information can be inferred from outputs, timing, or persistent state.

Recent systems explore how to make TEE-backed LLM inference practical under tight memory and performance constraints. \textbf{TZ-LLM} uses Arm TrustZone to protect on-device LLM execution by confining model weights and sensitive intermediate activations to the secure world and managing the host--TEE interface to reduce leakage and overhead~\cite{wang2025tz}.

\textbf{SPADA} and \textbf{PipeLLM} target secure, high-throughput inference by redesigning execution pipelines so protected computation can overlap with communication and scheduling; SPADA also explores scaling across multiple parties or nodes. These approaches aim to amortize TEE and cryptographic costs~\cite{chu2025spada,tan2025pipellm}. \textbf{Countermind} advocates defense in depth across the LLM-serving stack to address cross-layer attacks and side channels beyond the enclave boundary~\cite{schwarz2025countermind}.

Confidential Prompting studies privacy-preserving cloud inference when user prompts are sensitive. It proposes mechanisms to prevent the cloud provider from learning prompt contents while still returning model outputs~\cite{gim2024confidential}. This work complements TEE-based isolation by focusing on the prompt channel and motivates end-to-end threat models that account for possible leakage through execution traces or model-dependent side channels~\cite{gim2024confidential}.

TEE deployment also involves governance: Guitton et al. discuss attestation, model and data provenance, and operational responsibility as parts of end-to-end trust~\cite{guitton2024honest}.

\begin{tcolorbox}[colback=gray!8, colframe=gray!55, boxrule=0.4pt, left=1mm, right=1mm, top=0.8mm, bottom=0.8mm, title=Key findings]
% Hypervisor/TEE fuzzing is less mature than syscall/driver fuzzing due to limited instrumentation, weak failure signals (e.g., VM resets/enclave aborts), and fragmented vendor specifications. Still, tiny but highly privileged interfaces (hypercalls, virtual-device emulation, secure-world RPC) make bugs high impact. 
\revise{The examples discussed here include non-LLM hypervisor fuzzing, LLM-assisted TEE code annotation, and TEE-based protection of LLM inference.}
\revise{A promising future direction is to evaluate whether LLMs can help connect hardware/interface specifications to testable security properties of hypervisor and TEE implementations.}
% LLMs help infer argument schemas and setup sequences from headers/manuals/traces and, for non-crash violations, define semantic/differential oracles via replay across versions/configurations.
\end{tcolorbox}

\subsection{System Services} 
\label{sec:os:services}
% System services (e.g., \texttt{journald}, \texttt{sshd}), running in user space with elevated privileges, serve as the essential ``glue'' for system bootstrapping, configuration, and lifecycle management. While they share the security-critical nature of the kernel core discussed in \ref{sec:os:kernel}, these components orchestrate high-level OS behaviors and expose attack surfaces through complex runtimes and administrative daemons~\cite{kerrisk2010linux}. LLM-based techniques in this domain have transitioned toward semantic-aware auditing and enforcement, targeting the unique "User-Kernel bridge" through four specialized dimensions:
System services (e.g., \texttt{journald}, \texttt{sshd}), running in user space with elevated privileges, serve as the essential ``glue'' for system bootstrapping, configuration, and lifecycle management. While they share the security-critical nature of the kernel core discussed in \ref{sec:os:kernel}, these components orchestrate high-level OS behaviors and expose attack surfaces through complex runtimes and administrative daemons~\cite{kerrisk2010linux}. \smy{The following related techniques illustrate mechanisms potentially useful for service security; their evaluated targets differ:}
% (1) \textbf{semantic observability} to monitor system-level interactions using LLM-guided instrumentation~\cite{zheng2025agentsight},
\smy{(1) \textbf{semantic observability} of AI agents through eBPF-based correlation of LLM communications and system-level actions~\cite{zheng2025agentsight},}
(2) \textbf{intent-driven synthesis} of security policies and kernel extensions~\cite{gao2025offloading, zheng2024kgent},
% (3) \textbf{autonomous repair} of complex service logic and patches~\cite{xue2025pagent, mathai2025crashfixer},
\smy{(3) \textbf{autonomous repair} workflows, illustrated by post-processing of agent-generated patches in PAGENT and Linux-kernel repair in kAgent (formerly CrashFixer)~\cite{xue2025pagent,mathai2025crashfixer},}
% and (4) \textbf{formal enforcement} of runtime constraints via neural-symbolic oracles~\cite{wang2025agentspec, lyu2025ebpf}.
\smy{and (4) \textbf{runtime enforcement} of constraints on LLM agents in AgentSpec~\cite{wang2025agentspec}.}

\paragraph{Semantic Triage of Privileged Runtimes and Daemons}
Distinct from syscall-centric triage in Section~\ref{sec:os:kernel}, service analysis may involve varied runtimes, low-level daemons, stripped binaries, and heterogeneous log formats~\cite{xiong2023hext5,guan2024logllm}.

In log analysis, \textbf{LogPipe} constructs and updates a knowledge base of discrete and semantic log patterns. At inference time, it matches log sequences against this knowledge to guide the LLM, aiming to improve accuracy and explainability while reducing redundant queries~\cite{tao2026knowledge}. \textbf{LogLLM} uses the natural-language semantics of log messages to identify anomalous sequences, reducing reliance on explicit parsers and templates~\cite{guan2024logllm}.
\revise{The reported evaluations use general software-log datasets, not privileged OS service logs. LogPipe reports an average F1 of 97.5\% across eight public datasets; LogLLM evaluates on four public datasets}~\revcite{tao2026knowledge,guan2024logllm}.

At the binary-analysis layer, \textbf{HexT5} treats decompiled pseudo-code as a programming language and pre-trains a unified model for semantic inference. It evaluates tasks including binary-code summarization, identifier recovery, and similarity detection~\cite{xiong2023hext5}.

For differential testing of WebAssembly (WASM) runtimes, \textbf{LWDIFF} extracts and operationalizes knowledge from the Wasm specification to generate test cases, then compares runtime behavior to reveal semantic inconsistencies~\cite{zhou2025lwdiff}. \textbf{LLMSAN} requires LLM-driven bug detection to emit explicit data-flow paths and validates them using program-property decomposition and parsing-based checks~\cite{wang2024sanitizing}. \revise{Neither study evaluates privileged OS services}~\revcite{zhou2025lwdiff,wang2024sanitizing}.

\paragraph{Artifact Synthesis for Administration} 
\revise{Some adjacent work synthesizes network policies or other operational artifacts from natural-language intent.}

Service-related artifact synthesis translates natural-language goals into policies or configuration that govern interactions among OS components.

In network access control, \textbf{Clarify} introduces a disambiguation module that resolves semantic overlaps in ACL synthesis through active operator feedback, \revise{helping clarify the intended policy behavior}~\cite{mondal2025llm}.
\revise{In its evaluated network access-control setting, Clarify performs incremental synthesis: it generates a policy stanza and corresponding JSON specification, which a verifier checks before integration}~\revcite{mondal2025llm}.
\revise{A disambiguator asks targeted questions using differential examples to determine where the stanza belongs in the existing policy; the evaluation uses a small synthetic topology}~\revcite{mondal2025llm}.
% Beyond static configuration. 
\paragraph{Runtime Service Patching and Repairing}
 

Unlike the multi-step patch process used for kernel evolution in Section~\ref{sec:os:kernel}, service-repair work explores faster, fine-grained mitigations through self-healing operators~\cite{salva2025dynamic}. These approaches can generate supplementary code without modifying the original implementation and apply a fallback mitigation if the primary one fails~\cite{salva2025dynamic}.
\textbf{Salva et al.} use LLMs to generate encapsulation-based healing operators that implement such mitigations~\cite{salva2025dynamic}. \revise{Their evaluation targets RESTful applications, not privileged OS services}~\revcite{salva2025dynamic}.
% This methodology complements offline repair, such as \textbf{template-guided program repair}, which combines LLM generation with structured fix templates to ensure the long-term correctness of patches for system utilities\cite{huang2025template}. Empirical evaluations of code-trained LLMs demonstrate strong general code synthesis capability from natural-language specifications, motivating their use as patch generators within automated repair pipelines for system utilities\cite{chen2021evaluating}.
\revise{Template-guided repair combines LLM-generated candidates with structured fix templates in general software-repair benchmarks}~\revcite{huang2025template}.

\begin{tcolorbox}[colback=gray!8, colframe=gray!55, boxrule=0.4pt, left=1mm, right=1mm, top=0.8mm, bottom=0.8mm, title=Key findings]
\revise{The reviewed techniques address complementary parts of service auditing, analysis, and repair, often on different evaluation targets. Integrating them into a validated workflow for privileged services is a promising research direction.}
\end{tcolorbox}

\begin{comment}
\note{below are not related}
\paragraph{Agentic Sandboxing and Verification Infrastructure Hardening}
As system services increasingly adopt autonomous agentic architectures, security relies on enforcing safety bounds through the autoformalization of high-level requirements\cite{wang2025agentspec,gallardo2025agentic}.
While kernel verification targets low-level memory safety and isolation, service-level verification focuses on sandboxing agent interactions and hardening the underlying formal tools that provide security oracles\cite{wang2025agentspec}.

\textbf{AgentSpec} provides a framework for generating formal runtime monitors that sandbox LLM agents, ensuring their interactions with privileged system interfaces remain within safe boundaries\cite{wang2025agentspec}. 
Supporting this, \textbf{Agentic Lean Autoformalization (ALA)} translates natural language requirements into formal \textbf{LEAN} specifications to reduce the manual formalization burden\cite{gallardo2025agentic}.
\end{comment}
